\subsection{Admissibility Ratio $R_p$ Along the Chain $Q_8\to 2I\to\mathrm{LPS}_p\to SU(2)$}
  \label{subsec:Rp-along-chain}

  The admissibility ratio $R_p = A_{\rho_s}^{\max}/A_{\rho_1}^{\max}$ measures the
  relative advantage of the spinorial sector over the lowest non-trivial abelian
  sector.
  Its evolution along the chain of increasingly refined finite relational structures is:

  \begin{center}
    \begin{tabular}{lcc}
      \hline
      Structure & $R_p$ & Notes \\
      \hline
      $Q_8$ & $\sqrt{4/3}\approx 1.155$ & minimal non-abelian case \\
      $2I$ & $\sqrt{14/9}\approx 1.247$ & binary icosahedral \\
      $\mathrm{LPS}_5$ & $\approx 1.31$ & Ramanujan, $p=5$ \\
      $\mathrm{LPS}_{13}$ & $\approx 1.36$ & Ramanujan, $p=13$ \\
      $SU(2)$ (continuum) & $1$ & uniform spectrum, filter disappears \\
      \hline
    \end{tabular}
  \end{center}

  Several features of this table merit comment.

  First, $R_p$ is strictly greater than 1 for all finite structures in the chain.
  The spinorial sector always admits strictly more amplitude than the abelian sectors
  at finite resolution.

  Second, $R_p$ is non-monotone along the chain: it increases from $Q_8$ to the
  Lubotzky-Phillips-Sarnak family before converging to 1 in the continuum.
  The LPS graphs are Ramanujan ($|\text{non-trivial eigenvalues}|\leq 2\sqrt{d-1}$
  for $d$-regular graphs), meaning their spectral gap is maximal for their valency.
  This maximality does not itself maximize $R_p$; rather, it ensures that the
  admissibility filter is as selective as possible for a given graph degree, which
  is the relevant constraint for the continuum limit.

  Third, the convergence $R_p\to 1$ as the resolution increases reflects the fact
  that in $SU(2)\cong S^3$, all representation sectors coexist without hierarchy.
  The bounded-flux filter is non-trivial only at finite resolution; it disappears
  when the continuous group is restored.
